%% guide-p2a-src.tex　ガイドブック　保護者・先生のための本（大人用）の「出典」（組 G-B、頭文字 kgb）。親ファイルで付録（guide-p2b）のあとに読む。
%% \kgbsec・\kgbsub は guide-p2a.tex で定義（\providecommand）。材料：tools/guide-shidou.md 1節。
%% P1・P2 の URL のドメインは、材料では「co.jP」と書かれているが、ドメイン名は大文字・小文字を区別しないので「co.jp」と書いた。
\clearpage\kgbsec{}{出典}\label{kgb-src}
\begingroup
\urlstyle{same}\def\UrlFont{\sffamily\footnotesize}%
\Urlmuskip=0mu plus 1mu\relax
\expandafter\def\expandafter\UrlBreaks\expandafter{\UrlBreaks\do\a\do\b\do\c\do\d\do\e\do\f\do\g\do\h\do\i\do\j\do\k\do\l\do\m\do\n\do\o\do\p\do\q\do\r\do\s\do\t\do\u\do\v\do\w\do\x\do\y\do\z\do\0\do\1\do\2\do\3\do\4\do\5\do\6\do\7\do\8\do\9\do\_\do\-}%
この本で照らし合わせた資料です（2026年10月30日に調べました）。小学校・中学校・高等学校の指導要領の本文は、文部科学省の「学習指導要領コード」（告示本文を項目ごとに収めたデータ。M13）で、項目ごとに確かめました。小学校・中学校の本文は、教科書会社の新旧対照表（P1〜P3）と解説（M4・M5）とも照らしました。
\def\kgbsrc#1#2#3{\par\addvspace{3pt}\noindent\hangindent=3.2\zw\hangafter=1\makebox[3.2\zw][l]{\sffamily\bfseries\small #1}{\small #2}\par\noindent\hspace{3.2\zw}\parbox[t]{\dimexpr\linewidth-3.2\zw\relax}{\url{#3}}\par}%
\kgbsub{1}{文部科学省の資料（学習指導要領の本文・解説など）}
\kgbsrc{M1}{平成29・30・31年改訂学習指導要領（本文、解説）}{https://www.mext.go.jp/a_menu/shotou/new-cs/1384661.htm}
\kgbsrc{M2}{小学校学習指導要領（平成29年告示）}{https://www.mext.go.jp/content/1413522_001.pdf}
\kgbsrc{M3}{中学校学習指導要領（平成29年告示）}{https://www.mext.go.jp/content/1413522_002.pdf}
\kgbsrc{M4}{小学校学習指導要領（平成29年告示）解説 算数編}{https://www.mext.go.jp/content/20211102-mxt_kyoiku02-100002607_04.pdf}
\kgbsrc{M5}{中学校学習指導要領（平成29年告示）解説 数学編}{https://www.mext.go.jp/component/a_menu/education/micro_detail/__icsFiles/afieldfile/2019/03/18/1387018_004.pdf}
\kgbsrc{M6}{高等学校学習指導要領 比較対照表（数学）}{https://www.mext.go.jp/component/a_menu/education/micro_detail/__icsFiles/afieldfile/2018/07/13/1407085_5.pdf}
\kgbsrc{M7}{高等学校学習指導要領（平成30年告示）解説 数学編・理数編}{https://www.mext.go.jp/content/20230217-mxt_kyoiku02-100002620_05.pdf}
\kgbsrc{M8}{平成29年改訂の小・中学校学習指導要領に関するQ\&A（算数、数学に関すること）}{https://www.mext.go.jp/content/1422350_001.pdf}
\kgbsrc{M9}{中学校学習指導要領の特例を定める件（平成29年文部科学省告示第94号）}{https://www.mext.go.jp/a_menu/shotou/new-cs/__icsFiles/afieldfile/2017/07/11/1387780_003_1.pdf}
\kgbsrc{M10}{高等学校数学科教材「行列入門」（文部科学省）}{https://www.mext.go.jp/content/20230828-mxt-kyoiku01_000250597_1.pdf}
\kgbsrc{M11}{中央教育審議会資料「標準授業時数の在り方について」}{https://www.mext.go.jp/content/20200205-mxt_kyoiku01-000004708_9.pdf}
\kgbsrc{M12}{令和7年度大学入学共通テストにおける新教育課程対応について（大学入試センター提出資料）}{https://www.mext.go.jp/content/20211014-mxt_daigakuc02-000018441_3.pdf}
\kgbsrc{M13}{学習指導要領コード（文部科学省。小学校・中学校は平成29年告示、高等学校は平成30年告示の本文を項目ごとに表示したもの）}{https://jp-cos.github.io/}
\kgbsrc{M14}{大学入試センター「出題方法等の予告」（令和7年度大学入学共通テスト、数学）}{https://www.dnc.ac.jp/albums/abm.php?d=743&f=abm00003201.pdf}
\kgbsub{2}{教科書会社の新旧対照表（告示の本文を写したもの）}
\kgbsrc{P1}{大日本図書「小学校算数 現行・新学習指導要領 教授用資料 新旧対照表」}{https://www.dainippon-tosho.co.jp/sansu/files/cos2017e_sansu_f.pdf}
\kgbsrc{P2}{大日本図書「中学校数学 現行・新学習指導要領 教授用資料 新旧対照表」}{https://www.dainippon-tosho.co.jp/math/files/cos2017j_sugaku_f.pdf}
\kgbsrc{P3}{日本文教出版「小学校算数 新旧対照表」}{https://www.nichibun-g.co.jp/textbooks/sansu/download/new_cs/1705_newcs_sansu.pdf}
\kgbsub{3}{大学のシラバス・資料（8章）}
\kgbsrc{U1}{東北大学 工学部「工業数学Ⅱ」（2024）}{https://qsl.cds.tohoku.ac.jp/qsl/syllabus/display/tb44011}
\kgbsrc{U2}{東北大学 工学部「数学Ⅰ」（2024）}{https://qsl.cds.tohoku.ac.jp/qsl/syllabus/display/tb13011}
\kgbsrc{U3}{東北大学「線形代数学B」（2024）}{https://qsl.cds.tohoku.ac.jp/qsl/syllabus/display/cb42201}
\kgbsrc{U4}{京都工芸繊維大学 設計工学域「線形代数学Ⅱ」（2026）}{https://www.kit.ac.jp/wp/wp-content/uploads/sy2026/12061030.pdf}
\kgbsrc{U5}{北海道大学「微分積分学Ⅰ」（2026）}{https://www.math.sci.hokudai.ac.jp/~furuhata/ed/syllabus/26Calculus1.pdf}
\kgbsrc{U6}{神戸大学 工学部 前期シラバス}{https://www.office.kobe-u.ac.jp/eng-ofc/kym/syllabus/09_bz_a.pdf}
\kgbsrc{U8}{東京農工大学 工学部 数学科目一覧}{https://web.tuat.ac.jp/~mathsci/syllabus.html}
\kgbsrc{U9}{北海道大学「確率・統計入門」（2018）}{https://www.math.sci.hokudai.ac.jp/~furuhata/ed/syllabus/18IntroductionToStatistics.pdf}
\endgroup
